% Template for post or presentation abstracts submitted to ILAS 2022
% 1. Fill in the presenter's information (name, affiliation, email
% address) and talk information (title of presentation, abstract).
% 2. Compile to make sure there are no errors.
% 3. Save the .tex file for your abstract with a distinctive name,
% ideally "FamilynameGivenname.tex".
% 4. Upload your .tex file to https://clr.ie/129557
%
% * Please keep your LaTeX commands simple, and avoid the use of
%    user-defined macros.
% * Include details of co-authors and funding acknowledgements after the abstract.
% * Upload only the LaTeX file (with .tex extension).
% * Questions: Contact Niall Madden (Niall.Madden@NUIGalway.ie)

\documentclass[a4paper]{amsart}
\usepackage{amssymb}
\usepackage{hyperref}
\pagestyle{empty}
\date{}

%% IGNORE THE NEXT 4 LINES
\makeatletter %
\def\labelenumi{\theenumi.}
\def\theenumi{\@arabic\c@enumi}
\makeatother
%% STOP IGNORING NOW

\begin{document}

\title{QR decomposition of Chebyshev-Vandermonde matrices by means of VDR decomposition and the Arnoldi process}
\author{Ik-Pyo Kim} %Presenting author only
\address{Daegu University} % Affiliation of presenting author
\email{kimikpyo@daegu.ac.kr} % Email address of presenting author only

\maketitle

% The abstract  ***no more than one page***

 In this paper, we introduce a decomposition of Chebyshev-Vandermonde matrices, which is called the $VDR$ decomposition, where $V$ is an ordinary Vandermonde matrix, $D$ is diagonal and $R$ is upper triangular. Motivated by Brubeck et al. [Vandermonde with Arnoldi, SIAM Rev., 63 (2021), 405-415], we consider the Chebyshev-Vandermonde matrices with the $VDR$ decomposition and a $QR$ decomposition of $V$ such that the upper triangular matrix can be expressed recursively by means of the upper Hessenberg matrix in the Arnoldi process. By using the above results, we present an algorithm for solving linear systems with the Chebyshev-Vandermonde matrices as coefficient matrices. Moreover, we present an error analysis for the computed solutions by this algorithm and a generalized minimal residual algorithm, which is called GMRES.

%% Coauthor details here (optional - delete if not applicable)
% and funding acknowledgment (optional - delete if not applicable)
\emph{This is joint work with Arnold R. Kr\"auter (Montanuniversit\"at) and Seo-Ha Lee (Pusan National University).
Supported by This research was supported by Daegu University Research Grant 2018.}

\begin{thebibliography}{1}
\bibitem{my_key_for_OT49}
 Pablo D. Brubeck, Yuji Nakatsukasa, and Lloyd N. Trefethen.
\newblock Vandermonde with Arnoldi.
\newblock {\em  SIAM Rev.}  63(2):405–415,  (2021).

\bibitem{another_unique_key}
Youcef Saad and  Martin H. Schultz.
\newblock GMRES: {A} {Generalized Minimal Residual Algorithm for Solving Nonsymmetric Linear Systems}.
\newblock {\em  SIAM J. Sci. Comput.}, 7(3):856--869, 1986.
\end{thebibliography}


\end{document}


% Please leave the next bit alone
%%% Local Variables:
%%% mode: latex
%%% TeX-master: "../ILAS-2022"
%%% End:
